\newcommand{\DoNotLoadEpstopdf}{}
\documentclass[11pt,reqno]{amsart}
\usepackage[T1]{fontenc}
\usepackage{lmodern,amsmath,amssymb,amsthm,mathtools,mathrsfs,microtype}
\usepackage[margin=1.1in]{geometry}
\usepackage{xr-hyper}
\usepackage[hidelinks]{hyperref}
\externaldocument[supp-]{build/supplement}[supplement.pdf]
\hypersetup{pdftitle={Collision records, arithmetic, and raw local inversion in a triangular Lorentz gas},pdfauthor={Qian Qi}}
\setlength{\emergencystretch}{2em}
\allowdisplaybreaks[2]
\numberwithin{equation}{section}
\newtheorem{theorem}{Theorem}[section]
\newtheorem{lemma}[theorem]{Lemma}
\newtheorem{proposition}[theorem]{Proposition}
\newtheorem{corollary}[theorem]{Corollary}
\theoremstyle{definition}
\newtheorem{definition}[theorem]{Definition}
\newtheorem{inputstatement}[theorem]{Input statement}
\theoremstyle{remark}\newtheorem{remark}[theorem]{Remark}
\newcommand{\R}{\mathbb R}\newcommand{\Z}{\mathbb Z}
\newcommand{\one}{\mathbf1}\newcommand{\dd}{\,\mathrm d}
\newcommand{\esssup}{\operatorname*{ess\,sup}}
\newcommand{\diver}{\operatorname{div}}
\newcommand{\TV}{\operatorname{TV}}
\title[Collision records and raw local inversion]{Collision records, arithmetic, and raw local inversion in a triangular Lorentz gas}
\author{Qian Qi}
\date{October 12, 2026. A2-DYN, revision 89}
\subjclass[2020]{37D50, 37A50, 60F05, 28A75}
\keywords{Lorentz gas, exact collision records, critical Lorentz estimates, first-exit traces, microscopic concentration, raw local limits}
\begin{document}
\raggedbottom
\begin{abstract}
We study the pointwise local limit problem for the complete displacement,
collision-count, occupation and roof record of actual returns in a
triangular Lorentz gas. Retaining the measure of a tested roof set
through the complete residence sum yields a critical weak estimate
at exponent $32/31$ for the untruncated first-exit surface. Its
finite-count remainder remains summable down to residence zero.
The original raw density consequently has a uniform critical weak
bound. Together with the established local mass convergence, this
proves raw-error convergence at the critical exponent with every
positive logarithmic loss. The same logarithmic estimate acts directly
on the original critical-or-short-residence source. We also prove
convergence over the full central window in diffusive roof measure,
with the scale factor explicit. All exact labels, arithmetic classes,
physical weights and earlier proofs are retained. The organizing
pointwise assertion is unchanged; its residual essential-height
charge is not inferred to vanish from these distribution bounds.
\end{abstract}
\maketitle

\input{journal-problem-v89}

\subsection{Complete sources without a positive residence cutoff}
The pointwise error in \eqref{eq:v76-target} remains the organizing
objective. The new assertions estimate the same raw density in finite
local norms and the actual outgoing surface with all positive residence
lengths retained. These statements are not conditional on a microscopic
charge being zero.

For $1<q<\infty$ let $\mathcal N_q(f;J)$ be the measurable-set norm
in \eqref{eq:v88-weak-norm}, equivalent to weak $L^q$ on $J$.
For the original first-exit surface write $\overline{\mathcal F}_r$
for its average over exit widths between zero and $r$, as defined below.
\begin{theorem}[Untruncated surface and raw-density estimates]
\label{thm:v88-main}
Under the explicit inherited inputs of the next section, for every
$1<q<32/31$ and fixed $L>0$,
\[
 \limsup_m\sup_{R,n,k}\sup_{|J|=L}
 \mathcal N_q(m^2\overline{\mathcal F}_r;J)
                      \le C_{q,L}r^{1/q-31/32}.
\]
There is no lower gradient, residence or incidence cutoff in this
estimate. For every fixed $M,L>0$, the raw scalar error satisfies
\[
 \sup_{R,n,k}\sup_{|J|=L}
 \mathcal N_q\!\left(
 (m^2p_{n,R}(k,m,\cdot)-\mathcal L_{m,R}(k_1,k_2,\cdot,n))
 \one_{\mathcal C_{m,R,n,k}(M)};J\right)\longrightarrow0.
\]
The corresponding strong $L^p$ convergence holds for every
$1\le p<32/31$. All arithmetic classes remain unnormalized.
\end{theorem}
The proofs are Theorems~\ref{thm:v88-full-surface} and
\ref{thm:v88-raw-limit}. The new source theorem first establishes a
negative moment of total regular residence by a capped layer cake at
finite collision count. Its geometric trace application uses the
whole same-label segment; the source term is not replaced by a
boundary collar. A finite $q$ is essential to the conclusion. The
unchanged essential-supremum problem remains separate.

\subsection{Critical distribution and the original exceptional source}
Set $Q=32/31$ and $\gamma=1/32$. Theorem~\ref{thm:v89-endpoint-surface}
proves a finite-count measurable-set estimate for the complete surface
at the previously excluded exponent $Q$. Its proof keeps the factor
$\min\{|E|,\ell\}$ until after the residence sum; it does not
substitute an endpoint negative moment into a subcritical theorem.
Proposition~\ref{prop:v89-density-bound} then gives a critical weak
bound for the raw density itself.

The logarithmic norm in \eqref{eq:v89-log-norm} yields two further
statements. The raw error tends to zero in that norm for every
positive logarithmic exponent (Theorem~\ref{thm:v89-log-raw}).
The original weighted union of initial criticality and short total
residence has an explicit logarithmic modulus in the same critical
norm (Theorem~\ref{thm:v89-exceptional-log}). The latter estimate
acts on the exceptional source, not only its complementary surviving
fibers. Theorem~\ref{thm:v89-diffusive} gives the corresponding full
central-window conclusions with measure $dt/\sqrt m$.

These statements preserve \eqref{eq:v76-target} as the pointwise
objective. Neither a logarithmically weakened critical norm nor the
normalized central-window theorem is identified with essential height.
The exponent $32/31$ comes from the available residence exponent,
not a claim of physical optimality.

\input{journal-inputs-v89}

\input{core/205_capped_negative_residence_moments}
\input{core/206_uncut_surface_lorentz_control}
\input{core/207_exact_label_raw_density_local_norm}

\input{core/208_critical_surface_distribution}
\input{core/209_logarithmic_endpoint_and_exceptional_source}
\input{core/210_diffusive_window_and_exact_limit_order}

\section{The unchanged pointwise endpoint}
The new negative residence theorem is a finite-count assertion for an
actual singular insertion. Its full layer cake includes arbitrarily
short physical segments before the collision limit. The outgoing
trace theorem similarly sums all incident components before applying
a norm. Thus the complete surface weak-norm decay and raw local-norm
limit proved here require no positive residence cutoff.

These results do not identify a finite local norm with essential
height. The endpoint in \eqref{eq:v76-target} continues to require
\[
 \lim_{B\to\infty}\mathfrak R_{\varepsilon(B),M}=0,
\]
where $\mathfrak R$ is the original weighted-source concentration
charge of Supplement Theorem~\ref{supp-v87archive:thm:v87-endpoint}.
The critical distribution and logarithmic estimates constrain this same source but do not
exclude narrow peaks in its essential supremum. The older sufficient
pairs $(\mathfrak A,\mathfrak B)$ and $(\mathfrak A,\mathfrak Q)$
remain valid with their exact hypotheses in the complete archive.
No charge-zero assertion is inferred here.

In particular, a strong local $L^p$ theorem for $p<32/31$ is not a
uniform denominator theorem on positive arithmetic classes. The
scalar-first path remainder, collision-to-return clock transfer and
forward likelihood statements are preserved, not promoted to new
unconditional same-roof ratios. Zero arithmetic classes are never
normalized. All previous mathematical files and the complete preceding
article remain in the supplement and provenance.

\input{journal-references}
\end{document}
